%\chapter*{Glossary}
%================GLOSSARY =================%
% Terms that are not acronyms. The list is sorted alphabetically by name (or by the sort key).

\newglossaryentry{blackhole}{
	name={Blackhole},
	description={JMH object that consumes the value a benchmark produces, so that the JIT compiler
	cannot remove the measured code as dead code.}
}

\newglossaryentry{bootstrap}{
	name={Bootstrap (registry)},
	sort={bootstrap},
	description={Creation of a missing Function Registry file by listing the functions already
	deployed on the target provider (Cloud Run services on GCP, Lambda functions on AWS) and writing
	one entry per function found. It runs only when the deployed workflow has internal calls.}
}

\newglossaryentry{bucket}{
	name={Bucket},
	description={Storage container in Google Cloud Storage or Amazon S3. QuickFaaS uploads the
	deployment package to a bucket, from which the provider loads the function code; on AWS, the
	Lambda function must reside in the same region as its bucket.}
}

\newglossaryentry{cloudrun}{
	name={Cloud Run},
	description={Container-based serverless product of Google Cloud, also marketed as
	second-generation Cloud Functions. The Function Registry validates and discovers GCP
	functions through the Cloud Run API.}
}

\newglossaryentry{descriptor}{
	name={Deployment descriptor},
	description={JSON file read by QuickFaaS that holds a function's deployment parameters
	(provider, project, region, bucket, runtime, trigger) and the path of its code file. An internal
	call points to it through \texttt{deploymentDescriptorPath}; it is used only at Level~3 of the
	resolution cascade.}
}

\newglossaryentry{package}{
	name={Deployment package},
	description={Compressed source archive (ZIP) containing a function's code, its QuickFaaS
	template and the required libraries. It is the deployment format common to AWS, GCP and Azure,
	and the one QuickFaaS standardises on.}
}

\newglossaryentry{drift}{
	name={Drift},
	description={Difference between the endpoint stored in a registry entry and the current endpoint
	of the live function. Level~1 validation detects it and refreshes the entry.}
}

\newglossaryentry{endpoint}{
	name={Endpoint},
	description={Concrete address used to invoke a function: an HTTPS URL on GCP, or the function's
	ARN on AWS.}
}

\newglossaryentry{executionrole}{
	name={Execution role},
	description={IAM role that an AWS Lambda function assumes when it runs. When the deployment
	descriptor supplies none, OmniFlow creates one named
	\texttt{OmniFlow\allowbreak Lambda\allowbreak Execution\allowbreak Role}.}
}

\newglossaryentry{executionstep}{
	name={Execution step},
	description={Workflow step that performs concrete work, such as an HTTP call or a variable
	assignment (the Call and Assign constructs of the OmniFlow DSL).}
}

\newglossaryentry{externalcall}{
	name={External call},
	description={Call step that carries a literal \texttt{host}/\texttt{path}. It targets a function
	or API the developer does not own, bypasses the Function Registry and is never deployed through
	QuickFaaS.}
}

\newglossaryentry{firstgen}{
	name={First-generation Cloud Function},
	description={GCP function created through the first-generation Cloud Functions API, with a URL
	under \texttt{cloudfunctions.net}. It is what QuickFaaS deploys on GCP; its registry
	bindings are neither validated nor rediscovered, since both query only Cloud Run.}
}

\newglossaryentry{registry}{
	name={Function Registry},
	description={JSON file, one per provider, that maps each \texttt{functionRef} to the invocation
	metadata of an internal function. It is the persistence layer of the resolution cascade.}
}

\newglossaryentry{functionref}{
	name={Function reference},
	description={Stable logical identifier, written \texttt{functionRef}, that names an internal
	function in the workflow, the registry and the deployment descriptor. It is a name, not an
	endpoint; the form \texttt{"region/functionRef"} pins the lookup to a single region.}
}

\newglossaryentry{functionurl}{
	name={Function URL},
	description={HTTPS endpoint of an AWS Lambda function. The QuickFaaS AWS provider creates one,
	but OmniFlow stores and invokes the function's ARN instead.}
}

\newglossaryentry{hook}{
	name={Hook function},
	description={Cloud-agnostic function in which the developer writes the business logic, using
	the generic QuickFaaS request and response types instead of a provider SDK.}
}

\newglossaryentry{internalcall}{
	name={Internal call},
	description={Call step that declares its target with \texttt{internalFunction} instead of a
	literal endpoint. It targets a function the developer owns, whose endpoint is resolved at
	deployment time.}
}

\newglossaryentry{metadata}{
	name={Invocation metadata},
	description={Values stored in the Function Registry for each function: its
	\texttt{serviceName} and its \texttt{url} (an HTTPS URL on GCP, an ARN on AWS).}
}

\newglossaryentry{lockin}{
	name={Vendor lock-in},
	description={Dependence on one provider's SDKs, formats and services that makes moving to
	another provider costly. It appears at the function code level and at the workflow definition
	level.}
}

\newglossaryentry{warmup}{
	name={Warm-up},
	description={Benchmark iterations that run before the measured ones and whose results are
	discarded, so that the JIT compiler has already compiled the code when measurement starts. The
	evaluation uses three warm-up iterations of one second each.}
}

\newglossaryentry{orchestration}{
	name={Orchestration},
	description={Coordination of several functions by an external service that holds the control
	flow, the data flow and the error handling, instead of the functions themselves.}
}

\newglossaryentry{region}{
	name={Region},
	description={Geographic location where a provider runs its infrastructure (e.g.,
	\texttt{eu-west-1} or \texttt{europe-west1}). Every resource lives in a region, which is part of
	its identity.}
}

\newglossaryentry{rendering}{
	name={Rendering},
	description={Translation of OmniFlow's provider-agnostic workflow model into the
	provider-specific workflow definition, such as a Google Cloud Workflows YAML file or an AWS Step
	Functions state machine.}
}

\newglossaryentry{cascade}{
	name={Resolution cascade},
	description={Three-level procedure that binds each internal call to an endpoint at deployment
	time, stopping at the first level that succeeds: Level~1, a registry hit validated against the
	live function; Level~2, discovery on the provider; Level~3, deployment through QuickFaaS, only
	when the function is absent and a deployment descriptor is supplied.}
}

\newglossaryentry{roundrobin}{
	name={Round-robin},
	description={Distribution of items over a set of targets in turn, cycling back to the first
	after the last. In the benchmarks, the internal calls of a generated workflow reference the $F$
	distinct functions in turn.}
}

\newglossaryentry{serverless}{
	name={Serverless computing},
	description={Cloud execution model in which the provider allocates resources, patches the
	operating system and scales the infrastructure, and the user pays only for the time the code
	runs.}
}

\newglossaryentry{snapshot}{
	name={Snapshot (registry)},
	sort={snapshot},
	description={In-memory copy of the Function Registry that a resolver reads once per
	\texttt{resolve()} call and keeps in sync with every write, so that calls are resolved without
	re-reading the file.}
}

\newglossaryentry{stale}{
	name={Stale entry},
	description={Registry entry whose function no longer exists on the provider. It is removed and
	the deployment aborts, even when a deployment descriptor is present.}
}

\newglossaryentry{statemachine}{
	name={State machine},
	description={Workflow model in which each state performs work, makes a decision or controls the
	execution flow, passing data between states. AWS Step Functions represents workflows this
	way.}
}

\newglossaryentry{staticbinding}{
	name={Static binding},
	description={Embedding of the resolved endpoint in the rendered workflow at deployment time. The
	deployed workflow never consults the registry at run time, so an endpoint change takes effect
	only after the workflow is deployed again.}
}

\newglossaryentry{template}{
	name={Template function},
	description={Provider-specific wrapper maintained by QuickFaaS that is the entry point invoked
	by the provider. It adapts the provider's request to the generic format and calls the hook
	function.}
}

\newglossaryentry{workflow}{
	name={Workflow},
	description={Description of a process as a sequence of steps, with control flow and data flow,
	executed by a workflow service such as AWS Step Functions or Google Cloud Workflows.}
}

\newglossaryentry{workflowfirst}{
	name={Workflow-first deployment},
	description={Deployment approach in which the developer starts from a single workflow
	definition, and deploying it resolves, or if needed deploys, the functions it calls before
	deploying the workflow.}
}

% The text writes these terms as plain text, not with \gls, so add every entry to the list
% explicitly; glsignore drops the page number the list would otherwise show. As in acronyms.tex,
% this runs at the end of the document to avoid a blank page before the cover.
\AtEndDocument{\glsaddall[types={main},format=glsignore]}
